## Adversarial-review verdict

After reconstructing the proof independently rather than merely rerunning its audit scripts, I found:

\[
\boxed{\textbf{No fatal mathematical gap.}}
\]

My referee-style verdict is:

> **PROVISIONALLY VALID AFTER MINOR CORRECTIONS.**

That is stronger than “I could not disprove it,” but weaker than independent peer verification. I would **not retract the \(n/(12d^3)\) theorem on the evidence currently available**.

The original PDF is not literally valid as written because I found two genuine defects. Both are repairable without changing the theorem, proof strategy, or constant.

## The two real defects

### 1. Directed versus undirected boundary edges

The original manuscript mixes conventions when using the geodesic-load boundary estimate. The source argument counts a symmetric relation as \(nk\) **directed arcs**, while the manuscript sometimes calls the corresponding cut an undirected edge boundary. That creates an apparent factor-of-two problem.

I rederived the step from scratch. For a directed arc \(e\), let \(P_e\) be the total fractional load of all ordered-pair geodesics through \(e\). Coherent-configuration regularity makes this independent of \(e\), say \(P_e=P\). Then

\[
nkP=\sum_{x,y}\operatorname{dist}(x,y)\le dn^2,
\]

so

\[
P\le\frac{dn}{k}.
\]

For \(S\subset V\), every geodesic from \(x\in S\) to \(y\notin S\) uses an outward-directed crossing arc. Therefore

\[
|S|(n-|S|)
   \le |\partial^+(S)|P,
\]

and hence

\[
|\partial^+(S)|
\ge
|S|\frac{k}{d}\frac{n-|S|}{n}.
\]

Each undirected crossing edge corresponds to exactly one outward-directed crossing arc, so the numerical cut size is the same. **There is no lost factor of two.**

The published Pyber–Skresanov argument supplies the coherent-configuration geodesic uniformity and associated expansion machinery used here. citeturn238390view0

This is a mandatory notational correction, but not a mathematical failure.

### 2. Kivva theorem-number mismatch

Several citations use theorem numbers from Kivva’s arXiv version while the bibliography points to the published journal version. The content matches, but the numbering does not.

The relevant corrections include:

\[
\begin{aligned}
\text{arXiv Lemma 2.16} &\longrightarrow \text{journal Lemma 2.17},\\
\text{arXiv Lemma 2.17} &\longrightarrow \text{journal Lemma 2.18},\\
\text{arXiv Lemma 2.19} &\longrightarrow \text{journal Lemma 2.20},\\
\text{arXiv Theorem 2.24} &\longrightarrow \text{journal Theorem 2.25}.
\end{aligned}
\]

The journal version contains the needed geometric identities, local-clique description, local-eigenvalue theorem, multiplicity machinery, and Hamming/Doob endpoint classification. citeturn321420view0turn321420view1turn321420view2turn321420view3

Again, this is a source-interface defect rather than a flaw in the deduction.

## What survived the hostile reconstruction

### The adjacent-pair identity

I independently derived

\[
D(1)=2+\frac{2}{k}\sum_{i=2}^{d}k_i c_i.
\]

Because \(c_i\ge c_2=\mu\) for \(i\ge2\), and because \(c_2\le b_1\) implies \(k_2\ge k_1=k\), one obtains

\[
D(1)>\frac{\mu}{k}n.
\]

I found no missing endpoint or monotonicity assumption here.

### Small support implies strong local structure

For an automorphism with support density \(\rho\le1/2\), the repaired boundary argument gives a moved vertex \(x\) having at least

\[
\frac{k}{d}(1-\rho)
\]

fixed neighbors.

Under the manuscript’s small-support assumption, this exceeds \(\mu\), forcing

\[
x\sim x^g.
\]

Every fixed vertex is equidistant from \(x\) and \(x^g\), so the distinguishing set of the adjacent pair lies inside the support. Consequently,

\[
\mu<\rho k,
\qquad
\lambda>\frac{1-\rho}{d}k.
\]

This implication is sound.

### The full Metsch step and \(m\le d\)

I checked the ceiling term and all strict inequalities rather than using the usual coarse clique estimate. For

\[
\rho<\gamma=\frac1{12d^3},
\qquad
\alpha=\frac{1-\gamma}{d},
\]

the required scalar inequalities hold:

\[
\alpha^2>4\gamma,
\]

\[
\alpha-\frac{3\gamma}{2\alpha}>\frac1{d+1},
\]

and

\[
(d+1)^2\gamma<\alpha.
\]

The complete Metsch expression then produces a clique \(C\) satisfying

\[
|C|-1>\frac{k}{d+1}.
\]

Delsarte’s clique bound gives \(m<d+1\). Bang–Koolen’s geometricity criterion applies in the correct order, after which \(m\) is integral, yielding

\[
m\le d.
\]

I found no circular use of geometricity in this step.

### The geodesic Poincaré inequality

I reconstructed the proof with all orientations and factors of two exposed. It yields

\[
k-\theta_1
\ge
\frac{n^2k}
{\displaystyle\sum_{x,y}\operatorname{dist}(x,y)^2}
\ge\frac{k}{d^2}.
\]

For a mean-zero function \(f\), Cauchy–Schwarz along each geodesic gives

\[
2n\lVert f\rVert_2^2
\le
Q\sum_{e\ {\rm directed}}(\nabla_e f)^2
=
2Q f^{\mathsf T}(kI-A)f,
\]

while

\[
nkQ=\sum_{x,y}\operatorname{dist}(x,y)^2.
\]

The factors cancel exactly as claimed.

This lemma is not ornamental. The weaker published spectral-gap constant would not close the manuscript’s \(\mu=1\) branch at \(1/(12d^3)\).

### The difficult \(\mu=2\) branch

This received most of the audit.

I independently checked the exact recurrence

\[
y_{i+1}
=
\frac{k-\theta+c_i\,y_i/(1-y_i)}{b_i},
\qquad
y_i=\frac{u_{i-1}-u_i}{u_{i-1}}.
\]

I checked:

- the base case;
- the positivity range needed to divide by \(1-y_i\);
- the induction denominators;
- the product estimate for \(u_{t-1}\);
- the left and right distance-sphere tails;
- the concentration estimate \(k_t\ge(1-2m\varepsilon)n\);
- the complete multiplicity-certificate algebra;
- all branches in the \(R\)-factor argument;
- the case \(t=m<d\);
- the exact endpoint \(c_t=t=m=d\);
- the final deduction that every \(\psi_i=1\);
- exclusion of the Doob alternative by the valency inequality.

Outside the endpoint, the proof obtains

\[
k_{t-1}u_{t-1}^2
\ge
\frac nk\,FR
\]

with

\[
F>\frac{m}{m+1},
\qquad
R\ge\frac{m+1}{m}.
\]

Thus \(FR>1\), and Biggs’ multiplicity formula gives \(f_1<k\). The local-eigenvalue theorem then forces a local eigenvalue below \(-1\), contradicting the fact that the local graph is a disjoint union of cliques.

At the remaining endpoint, the intersection array is Hamming’s; Kivva’s cited endpoint classification leaves Hamming or Doob, and the manuscript’s large-valency estimate eliminates Doob. The relevant published Kivva results and their hypotheses are present in the cited source. citeturn321420view1turn321420view2turn321420view3

Most importantly, I found **no hidden reliance on the earlier computer-generated list of admissible tuples**. The current proof treats the parameter cases analytically.

## Two smaller clarifications I patched

The manuscript originally allowed the reader to infer \(y_2<1\) from a bound that is insufficient when \(m=2\). This is harmless because:

- when \(t=2\), the relevant product is empty and only \(u_1>0\) is needed;
- when \(t\ge3\), strict growth forces \(m\ge3\), and the displayed bound does give \(y_2<1\).

I made that split explicit.

At the Hamming endpoint, I also replaced a compressed sentence about the \(\psi_i\) with a maximal-index argument: assuming some \(\psi_j\ge2\), take the largest such \(j\), use \(\psi_{d-1}=1\), and derive the required contradiction. This eliminates an avoidable logical jump.

## Independent computation

I wrote a separate exact-rational checker from scratch rather than modifying the supplied one. It reconstructs the scalar constraints and every relaxed \(FR\)-case independently.

It passed through \(d=120\). The smallest relaxed certificate encountered was approximately

\[
FR=1.0079777384\ldots,
\]

still strictly above \(1\).

That finite test is only regression evidence. The manuscript’s infinite ranges are handled analytically—for example, the central \(F\)-estimate reduces to

\[
6d^3-10d^2-8d-5>0,
\]

and after writing \(d=x+3\), this becomes

\[
6x^3+44x^2+94x+43>0.
\]

## Bottom line

My current assessment is:

\[
\boxed{
\begin{array}{c}
\text{The original manuscript needs two mandatory corrections,}\\[2mm]
\text{but I found no fatal gap in the corrected proof of}\\[2mm]
\operatorname{motion}(X)\ge \dfrac{n}{12d^3}.
\end{array}
}
\]

The theorem would improve the published diameter dependence in this direction from \(d^{-6}\) to \(d^{-3}\), while retaining Johnson and Hamming as the exceptional families. Pyber and Skresanov’s published result supplies the \(d^{-6}\)-scale benchmark and the structural framework on which the manuscript builds. citeturn572917search0turn238390view0

I would now send the **patched** version, not the preceding PDF, to Claude and to a human specialist. The four best targets for an independent attack are:

1. uniformity of the directed geodesic loads in the coherent configuration;
2. the exact hypotheses and order of Metsch plus Bang–Koolen;
3. the passage \(FR>1\Rightarrow f_1<k\Rightarrow\) a forbidden local eigenvalue;
4. exhaustion of the \(R\)-factor cases and the Hamming/Doob endpoint.

Agreement between my audit and Claude’s would be meaningful evidence. It would still not replace review by someone working in distance-regular graphs or association schemes.

## Audit files

[Hostile referee report](sandbox:/mnt/data/babai_motion_adversarial_referee_report.pdf)

[Referee report source](sandbox:/mnt/data/babai_motion_adversarial_referee_report.tex)

[Adversarially patched proof](sandbox:/mnt/data/babai_motion_d3_adversarially_patched.pdf)

[Patched LaTeX source](sandbox:/mnt/data/babai_motion_d3_adversarially_patched.tex)

[Exact patch diff](sandbox:/mnt/data/babai_motion_d3_adversarial_patch.diff)

[Independent exact-rational checker](sandbox:/mnt/data/adversarial_independent_check.py)